\subsection{分次与符号}\label{sec:14d}
仍处于\stref{假设}{14.3}的情形，但另要求 $\dim(M^{\mathrm{fix}})=\dim(M)-2>0$，以免为平凡情形另作例外规定。假设 $M$ 带有二次复体积形式 $\eta_M^2$，不必 $\iota_M$ 不变。由于法丛平凡，它在不动点集上诱导同类形式；用局部平方根表示，就是
\[
\eta_M|M^{\mathrm{fix}}=\pm\pi^*(dz)\wedge\eta_{M^{\mathrm{fix}}}.
\]
令 $L$ 为满足\stref{假设}{14.4}且与不动点集相交的不变正合 Lagrangian 子流形。$L$ 上任意膜结构 $(\alpha^\#,P^\#)$ 都在不动点集
\SourcePage{215}
$L^{\mathrm{fix}}\subset M^{\mathrm{fix}}$ 上诱导膜结构。其分次为 $\alpha^{\mathrm{fix},\#}=\alpha^\#|L^{\mathrm{fix}}-\sigma/\pi$，其中 $\sigma$ 为法向角。\znote{207}{分次的这一改变量来自法向方向。按本书的约定，膜结构的分次记录的是从标架丛到 $\lambda^{\mathrm{top}}$ 平凡化的相位；而由 $L^{\mathrm{fix}}$ 的标架补上法向量 $e^{i\sigma}$ 才得到 $L$ 的标架，这一操作使 Maslov 势移动旋转角 $\sigma$（以弧度计），除以 $\pi$ 后即为 $\sigma/\pi$。因此只有当 $\sigma$ 是 $\pi$ 的整数倍时，$L^{\mathrm{fix}}$ 才原样继承 $L$ 的分次；一般情况下必须作此平移，引理 14.9 中的指标与定向空间才与 $L$ 上的相应对象一致。[来源：原书 §(14d)，印刷页 207。]}对于 Pin 结构，取与切丛关联的标架丛 $P\to L$ 和 $P^{\mathrm{fix}}\to L^{\mathrm{fix}}$，将 $L^{\mathrm{fix}}$ 的标架 $(X_1,\ldots,X_{n-1})$ 映为 $(e^{i\sigma},X_1,\ldots,X_{n-1})$，从而把后一个丛嵌入前一个；第一项是法向量。沿该嵌入拉回二重覆盖 $P^\#\to P$，并赋予继承的 $\operatorname{Pin}_{n-1}\subset\operatorname{Pin}_n$ 作用，就得到所需的 Pin 结构 $P^{\mathrm{fix},\#}$。

现在令 $(L_0^\#,L_1^\#)$ 为两个这样的带膜结构 Lagrangian 子流形，$(H,J)$ 为\stref{假设}{14.5}中的 Floer 数据。如前所述，将其限制到不动点集，便定义 $(L_0^{\mathrm{fix}},L_1^{\mathrm{fix}})$ 的 Floer 数据。取某个完全位于 $M^{\mathrm{fix}}$ 中的 $y\in\mathcal C(L_0,L_1)$，将同一对象视为 $\mathcal C(L_0^{\mathrm{fix}},L_1^{\mathrm{fix}})$ 的元素时记为 $y^{\mathrm{fix}}$。不动点集带有上述诱导膜结构时，有

\StatementNumber{9}\begin{lemma}\label{st:14.9}
$i(y^{\mathrm{fix}})=i(y)$ 且 $o(y^{\mathrm{fix}})\cong o(y)$。
\end{lemma}
\begin{proof}
按照\eref{11.39}，把 $i(y)$ 和 $o(y)$ 视为一对线性 Lagrangian 膜 $\Lambda_{k,y}^\#$ 的定向算子的指标与行列式线；这些线性膜来自 $L_k^\#$ 在 $y$ 处的切空间。我们断言，可以把定向算子选成 $\Z/2$ 等变的，使不变部分为 $y^{\mathrm{fix}}$ 的定向算子，反不变部分则可逆。这与\stref{引理}{14.6}和\stref{引理}{14.7}的证明类似。为简化记号，只在 $H=0$ 时写出细节。此时 $y$ 是 $L_0$ 与 $L_1$ 的不变交点，可如下构造定向算子 $D_{H,y}$。取向量丛 $E=H\times TM_y$，赋予非常值的复结构 $J_E$，其在无穷远处由 Floer 数据确定；也就是说，若 $\epsilon_\zeta$ 是 $H$ 唯一的条带状端，则 $s\ll0$ 时有 $J_E(\epsilon_\zeta(s,t))=J_t$。在 $E$ 上赋予平凡联络 $\nabla$。子丛 $F=F_\lambda\subset E|\partial H$ 由 $\operatorname{Gr}(TM_y)$ 中连接 $\lambda(0)$ 与 $\lambda(1)$ 的路径 $\lambda$ 给出。最后，要使其成为真正的定向算子，还须为 $\lambda$ 选择膜结构 $(\alpha_\lambda^\#,P_\lambda^\#)$；这当然限制了路径的同伦类。

在该同伦类内，可以选择 $J_E$ 在 $D\iota_M|TM_y$ 下不变，并使路径 $\lambda$ 由不变 Lagrangian 子空间组成，于是 $D_{H,y}=\bar\partial_\nabla$ 为 $\Z/2$ 等变。此外，$\lambda$ 的反不变部分，即 $\nu_y\cong\CC$ 中的一族直线 $\lambda^-(s)$，可取为 $\lambda^-(s)=e^{i\sigma(s)}\R$，其中函数 $\sigma:[0,1]\to\R$ 从 $\sigma(0)=\sigma_0$ 走到 $\sigma(1)=\sigma_1$。这样，不变部分按照上面处理 Lagrangian 子流形的同一程序继承诱导膜结构，使 $D^+_{H,y}$ 成为 $y^{\mathrm{fix}}$ 的定向算子。另一方面，条件\eref{14.10}在此就是 $\sigma_1-\sigma_0\in(-\pi,0)$，它蕴含 $\operatorname{index}(D^-_{H,y})=0$。再由\stref{引理}{11.5}，$D^-_{H,y}$ 可逆。
\end{proof}

\SourcePage{216}
\StatementNumber{10}\begin{note}\label{st:14.10}
也许值得用另一种方式重述此论证：避开指标理论，改用\eref{11.25}中指标和定向空间的初等定义。具体地，选择 $\lambda$，使其除上述条件外还属于 $\mathcal P^-\operatorname{Gr}(V)$。于是指标与行列式线只依赖于交 $\lambda(s)\cap\lambda(1)$。然而，可以使函数 $\sigma$ 严格递减；此时 $\sigma_1-\sigma_0\in(-\pi,0)$ 意味着对所有 $s<1$，均有 $\lambda^-(s)\cap\lambda^-(1)=0$。因此，这些交全部集中在不变部分。
\end{note}

假设 Floer 数据正则。取 $u\in\M_Z^*(y_0,y_1)$，它位于 $M^{\mathrm{fix}}$ 内，且为模空间的孤立点。如\stref{引理}{14.6}的证明，将同一映射视为 $\M_Z^*(y_0^{\mathrm{fix}},y_1^{\mathrm{fix}})$ 的元素时，记作 $u^{\mathrm{fix}}$。接下来断言：在\stref{引理}{14.9}的同构下，$u$ 与 $u^{\mathrm{fix}}$ 对 Floer 微分的贡献具有相同符号。用\secref{12f}的记号，就是下图交换：
\SourceEquationDiagram{diagram-14-13}{0.22015361}{14.13}
为看出这一点，须回顾\eref{12.2}：$c_u$ 是把线性化算子 $D_{Z,u}$ 与定向算子 $D_{H,y_1}$ 粘合，再将所得算子变形为 $D_{H,y_0}$ 得到的。在此，这些算子全是 $\Z/2$ 等变的；其不变部分可与 $M^{\mathrm{fix}}$ 中对应的算子认同，而后者给出 $c_{u^{\mathrm{fix}}}$；其反不变部分可逆，因而对行列式线的贡献平凡。此外，从 $D_{Z,u}\#D_{H,y_1}$ 到 $D_{H,y_0}$ 的同伦经过的都是同样形式的算子，故反不变部分始终可逆。同一论证也适用于\stref{引理}{14.7}情形中不变孤立伪全纯多边形的贡献 $c_{r,u}$。
